% ============================================================
%  Chapter 5 — Scenario Design
% ============================================================
\section{Scenario Design}
\label{ch:scenarios}

Six scenarios are defined to systematically explore the effect of
regulatory stringency, carbon pricing, and schedule flexibility on
optimal vessel operation. \Cref{tab:scenarios_overview} provides
a consolidated overview; the following sections describe each scenario's
rationale and parameter settings.

\begin{table}[h]
\centering
\caption{Scenario overview}
\label{tab:scenarios_overview}
\begin{tabular}{llllp{4.5cm}}
  \toprule
  ID & Name & Year & CII target & Key feature \\
  \midrule
  S1  & As-is validation     & 2022 & ---    & Replays observed speeds; no optimisation \\
  S2  & Max speed            & 2026 & ---    & All arcs at 16~kn; upper-bound emissions \\
  S3  & CII-optimal 2026     & 2026 & Rating C & Min-cost under 2026C threshold; fixed schedule \\
  S4  & Strict 2030 + ETS    & 2030 & Rating B & 2030B threshold + EUR 100/t\,CO$_2$; \textbf{infeasible} \\
  S4b & Relaxed 2030 + ETS   & 2030 & Rating C & As S4 but Rating C threshold; feasible \\
  S5  & Flexible schedule    & 2026 & Rating C & As S3 but global cycle-time cap replaces per-arc budgets \\
  \bottomrule
\end{tabular}
\end{table}

% ── 5.1 S1 — As-is validation ─────────────────────────────────
\subsection{S1: As-Is Validation}
\label{sec:scen:s1}

S1 replays the mean observed speed on each arc from the EU \MRV{}
dataset. No optimisation is performed; speeds are rounded to the nearest
integer knot and clamped to the model speed range $[10,16]$~kn. S1
serves two purposes: it validates the fuel model against reported
emissions (\cref{sec:case:validation}), and it establishes a
benchmark representing the vessel's historical operation.

% ── 5.2 S2 — Max speed ────────────────────────────────────────
\subsection{S2: Max Speed}
\label{sec:scen:s2}

S2 forces 16~kn on every arc, representing the worst-case emissions
scenario and providing an upper bound on fuel consumption and \CII{}. It
is not subject to any \CII{} constraint and requires no optimisation.
Comparison with S3 quantifies the cost of \CII{} compliance under the
2026 regulation.

% ── 5.3 S3 — CII-optimal 2026 ─────────────────────────────────
\subsection{S3: CII-Optimal 2026}
\label{sec:scen:s3}

S3 minimises total voyage cost subject to the 2026 Rating~C \CII{}
threshold (15.54~\gGTnm{}) and the published per-arc time budgets. No
\ETS{} carbon cost is included. This scenario models the minimum-cost
operation the vessel must adopt to maintain compliance with the current
regulatory trajectory.

\paragraph{Parameter values.}
\begin{itemize}
  \item Fuel price: EUR~800/t (2026 indicative HFO value).
  \item CO$_2$ price: EUR~40/t (EU \ETS{} at full 100\,\% coverage from 2026,
        assumed EUA price EUR~40/t in the near-term trajectory).
  \item \CII{} threshold: 15.54~\gGTnm{} (Rating C, 2026).
\end{itemize}

% ── 5.4 S4 — Strict 2030 + ETS ───────────────────────────────
\subsection{S4: Strict 2030 + ETS}
\label{sec:scen:s4}

S4 combines a 2030 Rating~B \CII{} threshold (9.93~\gGTnm{}, assuming
a 31\% reduction factor) with EU \ETS{} carbon pricing at
EUR~100/t\,CO$_2$. This scenario represents a stringent
\textit{aspirational} 2030 regulatory environment.

S4 is \textbf{infeasible}: the published schedule forces minimum
speeds on four arcs (\cref{tab:bottleneck_arcs}) that collectively
produce an aggregate attained \CII{} of 10.32~\gGTnm{} --- above
the 9.93~\gGTnm{} target. The infeasibility is structural: no
combination of speed choices satisfying the time budgets can achieve
Rating~B. This finding is reported as a thesis contribution.

The minimum achievable attained \CII{} under the published schedule is
computed analytically by the \texttt{\_compute\_min\_achievable\_cii()}
function: for each arc, the lowest integer-knot speed satisfying the time
budget is selected, fuel consumption is evaluated at that speed, and the
resulting aggregate CO$_2$ is divided by the effective capacity--distance
product.  The function returns 10.3187~\gGTnm{} for the base instance
(22 arcs, $n \geq 2$ filter, mean cargo), confirming the infeasibility
of S4 without any reliance on Gurobi.

% ── 5.5 S4b — Relaxed 2030 + ETS ────────────────────────────
\subsection{S4b: Relaxed 2030 + ETS}
\label{sec:scen:s4b}

S4b retains the 2030 economic conditions (EU \ETS{} at EUR~100/t,
fuel price EUR~914/t to reflect 2030 market estimates) but relaxes the
\CII{} target to Rating~C (12.05~\gGTnm{}). This scenario is feasible
with 14.4\% slack and represents the cost-optimal 2030 operation
when the vessel settles for Rating~C rather than Rating~B.

\paragraph{Regulatory note.}
Post-2026 reduction factors have not been formally adopted at the time
of writing. The 31\% reduction assumed here is a scenario assumption
based on the \IMO{} 2023 GHG Strategy's 40\% fleet-wide ambition
\parencite{mepc377} and the confirmed 11\% reduction for 2026.
Actual post-2026 factors may differ.

% ── 5.6 S5 — Flexible schedule ───────────────────────────────
\subsection{S5: Flexible Schedule}
\label{sec:scen:s5}

S5 relaxes the per-arc time budgets and replaces them with a single
global cycle-time cap $T_{\max} = 1{,}240$~h. All other parameters
are identical to S3. This scenario quantifies the operational and
economic value of schedule flexibility: by allowing the vessel to
redistribute sailing time across legs, it can operate at lower speeds
on tight-budget arcs while still completing the cycle on time.

\paragraph{Comparison with S3.}
The difference between S5 and S3 directly quantifies the \textit{value
of schedule flexibility}: the additional fuel and cost savings achievable
if the operator negotiates relaxed port windows, compared to operating
within the published fixed timetable.

% ── 5.7 Scenario parameter summary ───────────────────────────
\subsection{Scenario Parameter Summary}
\label{sec:scen:params}

\begin{table}[h]
\centering
\caption{Scenario economic and regulatory parameters}
\label{tab:scen_params}
\begin{tabular}{lrrrlr}
  \toprule
  ID & {Fuel (EUR/t)} & {CO$_2$ (EUR/t)} & {\CII{} threshold} & Schedule & {Cycle cap (h)} \\
  \midrule
  S1  & 800  & 0   & ---     & Observed speeds & --- \\
  S2  & 800  & 0   & ---     & 16 kn fixed     & --- \\
  S3  & 800  & 40  & 15.5412 & Per-arc budget  & --- \\
  S4  & 914  & 100 & 9.9291  & Per-arc budget  & --- \\
  S4b & 914  & 100 & 12.0488 & Per-arc budget  & --- \\
  S5  & 800  & 40  & 15.5412 & Global cap      & 1240 \\
  \bottomrule
\end{tabular}
\end{table}
